\subsection{Problema del Logaritmo discreto}
Il problema del logaritmo discreto è alla base di molti sistemi di cifratura e protocolli di crittografia.


\begin{definition}[Logaritmo Discreto]
    Per un gruppo moltiplicativo $(G,\cdot)$ e un elemento $\alpha \in G$ di ordine $n$ e sia $\beta \in \langle \alpha \rangle$ trovare un intero $a$ tale che $0 \leq a \leq n-1$ e
    \[
        \alpha^{a} = \beta
    \]
    ovvero ($a = log_{\alpha}(\beta)$).
\end{definition}

\vspace{1em}p 
La sua particolare utilità è fondata dalla difficoltà di invertire l'operazione di esponenziazione.

\subsection{Crittosistema di ElGamal}
Il sistema crittografico di ElGamal è basato sul problema del logartimo discreto nel gruppo $(\mathbb{Z}_p^*,\cdot)$.

\begin{definition}[El Gamal]
    Sia $p$ un numero primo tale che il logaritmo discreto nel gruppo $(\mathbb{Z}_p^*,\cdot)$ è intrattabile, sia $\alpha \in \mathbb{Z}_p^*$ un elemento primitivo.
    Il crittosistema ha $\mathcal{M} = \mathbb{Z}_p^*,\  \mathcal{C} = (\mathbb{Z}_p^* \times \mathbb{Z}_p^*)$, infine $\mathcal{K}$ è definito come segue
    \[
        \mathcal{K} = \{ (p,\alpha,a,\beta) \mid \beta \equiv \alpha^a \pmod p\}
    \]
    le variabili $p,\alpha,\beta$ formano la chiave pubblica e $a$ rappresenta la chiave privata
\end{definition}